% Teil II: der durchgerechnete Weg.

\subsection*{Gegeben und gesucht}
$a_1=1$, $a_{n+1}=\frac{1+a_n}{2+a_n}$. (a) $a_2,\dots,a_5$; (b) Beweis
$a_n^2+a_n>1$ für alle $n$; (c) monoton fallend, Konvergenz, Grenzwert.

\subsection*{a) Anfangswerte}
\begin{align}
  a_2&=\frac{1+1}{2+1}=\frac23, &
  a_3&=\frac{1+\frac23}{2+\frac23}=\frac{5/3}{8/3}=\frac58,\\
  a_4&=\frac{1+\frac58}{2+\frac58}=\frac{13/8}{21/8}=\frac{13}{21}, &
  a_5&=\frac{1+\frac{13}{21}}{2+\frac{13}{21}}=\frac{34/21}{55/21}=\frac{34}{55}.
\end{align}

\subsection*{b) Beweis $a_n^2+a_n>1$}
Induktion. IA ($n=1$): $1^2+1=2>1$. \checkmark
IS: Angenommen $a_n^2+a_n>1$. Mit der Rekursion multipliziert man um:
\begin{align}
  a_{n+1}^2+a_{n+1}-1
  &=\frac{(1+a_n)^2}{(2+a_n)^2}+\frac{1+a_n}{2+a_n}-1 \\
  &=\frac{(1+a_n)^2+(1+a_n)(2+a_n)-(2+a_n)^2}{(2+a_n)^2} \\
  &=\frac{a_n^2+a_n-1}{(2+a_n)^2}.
\end{align}
Nenner $>0$ und Zähler $a_n^2+a_n-1>0$ (Induktionsannahme), also
$a_{n+1}^2+a_{n+1}-1>0$, d.h. die Behauptung für $n+1$. Aus IA und IS folgt
$ a_n^2+a_n>1$ für alle $n$.

\subsection*{c) Monotonie und Grenzwert}
\begin{equation}
  a_{n+1}-a_n=\frac{1+a_n}{2+a_n}-a_n
  =\frac{1-a_n-a_n^2}{2+a_n}.
\end{equation}
Wegen $a_n^2+a_n>1$ ist der Zähler negativ, $2+a_n>0$; also
$a_{n+1}<a_n$: streng monoton fallend. Mit $a_n>0$ (untere Schranke) ist
$(a_n)$ nach unten beschränkt, also konvergent gegen $L$. Fixpunkt:
\begin{equation}
  L=\frac{1+L}{2+L}\ \Longrightarrow\ L^2+L-1=0\ \Longrightarrow\
  L=\frac{-1\pm\sqrt5}2 .
\end{equation}
Alle $a_n>0$ und fallend mit $L\ge0$, also
\begin{equation}
  L=\frac{\sqrt5-1}2 .
\end{equation}

\begin{ergebnis}
\begin{equation*}
  \text{a)}\ \tfrac23,\tfrac58,\tfrac{13}{21},\tfrac{34}{55},\qquad
  \text{b)}\ a_n^2+a_n>1,\qquad
  \text{c)}\ \text{monoton fallend},\ L=\frac{\sqrt5-1}2 .
\end{equation*}
\end{ergebnis}

\begin{kontrolle}
b) mit $a_2=\frac23$: $\frac49+\frac23=\frac{10}9>1$; c) Grenzwert
$\frac{\sqrt5-1}2\approx0.618$, die Folge fällt von $1$ zu diesem Wert --
passt zu $a_{n+1}<a_n$ und $L>0$.
\end{kontrolle}

\begin{merksatz}
Bei rekursiven Folgen: Glieder mit Brüchen ausrechnen, Ungleichungen per
Induktion, Monotonie über die Differenz, Grenzwert über die
Fixpunktgleichung -- die positive Lösung nehmen.
\end{merksatz}
